\documentclass[12pt]{article}
\usepackage[T1]{fontenc}
\usepackage[utf8]{inputenc}
\usepackage[brazil]{babel}
\usepackage[margin=1in]{geometry}
\setlength{\parindent}{0pt}
\begin{document}
\section*{Tema 73}
Processo: PEDILEF 2006.63.01.052381-5/ SP\\
Situacao: Julgado\\
Ramo: DIREITO PREVIDENCIÁRIO\\
Questao: Saber qual a composição do grupo familiar para concessão do benefício assistencial, no período anterior à Lei n. 12.453/2011.\\
Tese: O grupo familiar deve ser definido a partir da interpretação restrita do disposto no art. 16 da Lei n. 8.213/91 e no art. 20 da Lei n. 8.742/93, esta última na sua redação original.\\
Fonte: https://www.cjf.jus.br/phpdoc/virtus/
\end{document}
